\documentclass[10pt,a4paper]{amsart}
\usepackage[margin=19mm]{geometry}
\usepackage{amsmath,amssymb,mathtools}
\usepackage[scr=rsfs,cal=euler]{mathalfa}
\usepackage{xfrac}
\usepackage[unicode,hidelinks,bookmarks=false]{hyperref}
\newcommand{\ie}{i.e.\ }
\newcommand{\cf}{cf.\ }
\newcommand{\resp}{respectively\ }
\newcommand{\Subsection}{\subsection}
\providecommand{\nobreakdash}{\nobreak\hbox{-}}
\setlength{\emergencystretch}{2em}
\multlinegap0pt
\usepackage{tikz}
\usepackage{bm}
\usepackage{amssymb}
\usepackage{enumerate}
\newtheorem{proposition}{Proposition}
\title{Hausdorff measure of the free boundary for the $p$-obstacle problem with subcritical exponents}
\author{Jing Yu, Jun Zheng}
\date{Source: arXiv:2603.23121v1 (CC BY 4.0)}
\begin{document}
\maketitle

\noindent\emph{Source and licence.} Hausdorff measure of the free boundary for the $p$-obstacle problem with subcritical exponents. Version arXiv:2603.23121v1 (CC BY 4.0), distributed by arXiv under the Creative Commons Attribution 4.0 International licence, http://creativecommons.org/licenses/by/4.0/, as recorded in the arXiv licence metadata for this exact version and shown on the abstract page https://arxiv.org/abs/2603.23121v1. Full licence notice: \texttt{licenses/arxiv-2603.23121-CC-BY-4.0.txt}.\par\smallskip

\noindent\textit{Editorial scope.} Counted unit: the article's proposition nondegeneracy (source0.tex lines 1045--1051). The article's complete original proof of it is delivered with it (source0.tex lines 1053--1115, one \texttt{proof} environment of 63 lines). No mathematics is written, repaired or re-derived: the statement and the proof are the article's own lines, in the article's own notation, and no retained line is substituted or re-flowed.  Retained uncounted as the source-local context the proof needs: lines 1029--1037.  Nothing was omitted from any retained span, so the package carries no omission record.  No retained line contains a citation, so the wrapper carries no bibliography and no bibliography entry.  No retained line contains a figure environment, an image-inclusion command or drawing code, so no image is delivered and the package declares no asset dependency.  Editorial formatting only, in addition: an AMS \texttt{amsart} preamble that restates the article's own declarations and macros used by the retained lines, namely the environments \texttt{proposition} on separate counters, together with the standard packages those lines need.  The retained environments print in this document as Proposition 1 in order of appearance, whereas the article prints the same items with the numbering of its own full document.  The complete native source of the article as distributed in the arXiv e-print for this exact version is bundled unchanged as \texttt{sources/197055/source0.tex}.
	\begin{proposition}\label{grow}
		Let $y \in \Upsilon^+$ and $B_{r_1}(y) \subset \subset \Omega$  for some $r_1 > 0$. Then there exists a positive constant $C_1$ depending only on $N$, $p$, $\lambda$, $a_0$, and $a_1$ such that
		\begin{align*}
			|u(x)| \leq C_1 |x - y|^{\frac{p}{p-1}},\forall x \in B_r(y) \quad\text{and}\quad |\nabla u(x)| \leq C_1 |x - y|^{\frac{1}{p-1}},  { \forall x \in B_r(y) }
%\notag\\
			%|\nabla u(x)| \leq C_1 |x - y|^{\frac{1}{p-1}}, { \forall x \in B_r(y),}\notag
		\end{align*}
		hold  for every $r \in (0, r_1)$.
	\end{proposition}

	\begin{proposition}\label{nondegeneracy}
		Let $y \in \Upsilon^+$ and let $B_{r_2} \subset\subset \Omega$ with $	r_2 \leq \min\left\{r_1, \frac{1}{a_1} , \left(\frac{m_1}{2m_2 C_1^{\lambda-1}}\right)^{\frac{p-1}{p(\lambda-1)}}\right\}$. Then there exists a positive constant $C_2$ depending only on $N$, $p$, $m_1$, $a_1$, and $r_2$ such that
		\begin{align}
			\sup_{\partial B_r(y) \cap \{u > 0\}} u \geq C_2 r^{\frac{p}{p-1}},\forall r \in \left( 0, \frac{r_2}{2}\right) ,\notag
		\end{align}
		where $r_1$ and $C_1$ are positive constants the same as in Proposition~\ref{grow}.
	\end{proposition}

	\begin{proof}
		Let $y\in\{u>0\}$. For any $x_0\in\Upsilon^+$ satisfying $|y-x_{0}|\leq\frac{r_2}{2}$, we have $B_{r}(y)\subset B_{r_2}(x_{0})$ with $r<\frac{r_2}{2}$. Define $C_2:=\frac{p-1}{p}\left(\frac{m_1}{2N(a_1+1)}\right)^{\frac{1}{p-1}}$ and $v(x):=C_2|x-y|^{\frac{p}{p-1}}$ in $B_{r}(y)$. By  direct computations, we have
		\begin{align}
			\nabla v&=C_2\frac{p}{p-1}|x-y|^{\frac{2-p}{p-1}}(x-y),\notag\\
			|\nabla v|&=C_2\frac{p}{p-1}|x-y|^{\frac{1}{p-1}},\notag\\
			|\nabla v|^{p-2}&=C_2^{p-2}\left( \frac{p}{p-1}\right) ^{p-2}|x-y|^{\frac{p-2}{p-1}}\notag.
		\end{align}
		
		Noting that $r_2 < \frac{1}{a_1}$ and using the definition of $C_2$,  we obtain
		\begin{align}
			%\Delta_p v &=
			\mathrm{div}\left(a(x)|\nabla v|^{p-2}\nabla v\right)
			=
			&\mathrm{div}\left(C_2^{p-1}\left(\frac{p}{p-1}\right)^{p-1}a(x)(x-y)\right)
			\notag\\ =
			&C_2^{p-1}\left(\frac{p}{p-1}\right)^{p-1}\mathrm{div}\left( a(x)(x-y)\right)
			\notag\\ =
			&NC_2^{p-1}\left(\frac{p}{p-1}\right)^{p-1}a(x) + C_2^{p-1}\left(\frac{p}{p-1}\right)^{p-1}(x-y)\nabla a(x)
			\notag\\\leq
			& Na_1C_2^{p-1}\left(\frac{p}{p-1}\right)^{p-1} + Na_1r_2C_2^{p-1}\left(\frac{p}{p-1}\right)^{p-1}
			\notag\\ \leq
			&Na_1C_2^{p-1}\left(\frac{p}{p-1}\right)^{p-1} + NC_2^{p-1}\left(\frac{p}{p-1}\right)^{p-1}
			\notag\\ =
			&N(a_1+1)C_2^{p-1}\left(\frac{p}{p-1}\right)^{p-1}
			\notag\\ =
			&\frac{m_1}{2}\   \text{in}\   B_{r}(y).\label{6.1}
		\end{align}
		Using the fact that  $B_r(y) \subset B_{r_2}(x_0)$ and Proposition~\ref{grow}, we deduce that
		\begin{align*}
			\mathrm{div}\left(a(x)|\nabla u|^{p-2}\nabla u\right)
			=
			&  m_1 - m_2 u^{\lambda-1}
			\notag\\ \geq
			&  m_1 -  m_2 C_1^{\lambda-1} r_{2}^{\frac{p(\lambda-1)}{p-1}}
			\notag\\ \geq
			&\frac{m_1}{2}\   \text{in}\   B_r(y) \cap \{u > 0\},%\label{6.3}
		\end{align*}
		which, along with \eqref{6.1}, implies that
		\begin{align}
			-\mathrm{div}\left(a(x)|\nabla u|^{p-2}\nabla u\right) \leq -\mathrm{div}\left(a(x)|\nabla v|^{p-2}\nabla v\right)\   \text{in}\   B_r(y) \cap \{u > 0\}.\notag
		\end{align}
		Note that the definitions of $u$ and $v$ ensure that
		\begin{align}
			u=0\leq v \text{ on } B_{r}(y)\cap\Upsilon^{+}.\notag
		\end{align}
		If $u\leq v$ on $\partial B_{r}(y)\cap\{u>0\}$, then by the comparison principle, we obtain
		\begin{align}
			u\leq v  \text{ in } B_{r}(y)\cap\{u>0\}.\notag
		\end{align}
		However, $u(y)>0=v(y)$, which yields a contradiction. Therefore, there exists $y_{0}\in\partial B_{r}(y)\cap\{u>0\}$ such that
		\begin{align}
			u(y_{0})\geq v(y_{0})=C_2 r^{\frac{p}{p-1}},\notag
		\end{align}
		so we have
		\begin{align}\label{6.7}
			\sup_{\partial B_{r}(y)\cap\{u>0\}}u\geq C_2r^{\frac{p}{p-1}}.
		\end{align}
		Now for $y \in \partial\{u > 0\}$, we take a sequence $\{y_i\} \subset \{u > 0\}$ such that $y_i \to y$ as $i \to \infty$. By \eqref{6.7}, we have
		\begin{align}
			\sup_{\partial B_r(y_i) \cap \{u > 0\}} u \geq C_2 r^{\frac{p}{p-1}},\notag
		\end{align}
		which, along with the continuity of $u$, implies the desired result.
	\end{proof}

\end{document}
